% LNCS: abstract 150--250 words. No p-values in the abstract. Numbers: facts.md (Test-4: 40, 6, 2 families; ΔU −3.1;
% H1 not significant; tokens and latency of v3 against v1), adaptive rounds, AgentDojo v3.
% 15-page version (2026-11-20), condensed from the long version in commit 6814b3a.
% 20/11: synchronised with the Vietnamese version after owner feedback (round 1 and language pass).
% 20/11 (round 3b, English only): wording and connectives pass, Figure 2 on with [H].
% 20/11 (round 3c): notation only where equations need it, the request in words elsewhere, gate conditions as a list, plainer claims.
\begin{abstract}
Instructions planted in content that an LLM agent reads through its tools can make the agent send data to an unknown
destination, leak confidential values or alter a legitimate action. This paper proposes \ours, a deterministic
policy-enforcement mechanism at the Tool Broker that combines three components without any extra LLM call: an intent
policy that is extracted from the user request and handles negation, sensitivity labels for individual values, and
pinning of the values the user stated. Authorisation safety, confidentiality and integrity are proved on a
state-transition model under five explicit assumptions. On a Vietnamese test set of 48 attack families, sealed before
any code was written, \ours{} reduces the number of harmed families on three open-weight backbones from 40 with no
defence to 2. With string-matching rules in place of the three components, 6 families remain harmed. \ours{} blocks four
of these, each with the component designed for it, lets no new family through, and shows the same direction on every
backbone. Because string matching is already near the floor, confirming this difference with the pre-registered test
requires a larger test set. Utility drops by 3.1 percentage points relative to no defence, and tokens and latency
stay within 1.4\% of string matching. In two rounds of adaptive attacks, no family succeeds that does not already
succeed with the original attack. On AgentDojo in English and Vietnamese, attack success stays low, at the cost of
utility on tasks whose destinations come from tool data.

\keywords{Indirect prompt injection \and LLM agents \and Intent-based authorisation \and Information flow control \and
Vietnamese}
\end{abstract}
